\chapter{Vulnerability Analysis and Hardening}
\label{chap:vulnerability-analysis}

This chapter relates each observed original-system weakness to its code-level analysis, implemented mitigation, and available post-implementation evidence. For every attack A01--A08, a descriptive hardened-package evaluation is followed by a matched single-control ablation. The superseding A04 and earlier matched A06 ablations use family-label-omitted prompts, while the primary A07 ablation evaluates only the natural A07-N family. A later prospectively frozen A06 challenge separately creates a behaviourally active guards-off canary condition, and a supplemental matched A07-S experiment restores the historical synthetic-trigger family. Each result applies only to its tested prompt condition. Family-label omission denotes removal of the explicit family label; it does not imply semantic neutrality.

\subsection*{Descriptive package comparison}

The descriptive baseline in Chapter~\ref{chap:baseline-results} reports the failures observed in the full attack-family-label-free execution of the original implementation. The original hardened-package matrix reports what was observed after the set of hardening changes. Their comparison provides descriptive package evidence. Incomplete source-state manifests and attack-specific prompt provenance prevent an isolated causal estimate.

A01 target-level prompts are partially preserved/reconstructable but not package-equivalent. A02 complete user-prompt sequences are preserved, but their final attack prompts differ systematically. A03 final attack prompts match, but its protected seed prompts---the authorised setup requests used to place the target value in conversational state before the access downgrade---differ. A05 final attack prompts also differ. Exact prompt text is unavailable for the historical A04 and A08 package runs.

The later superseding A04 matched ablation has preserved prompt sequences that are identical across the off/on arms. A06 attack and warm-up user prompts match at the stored record level. Its package comparison spans two implementation stages and therefore does not isolate one hardening component.

Figure~\ref{fig:package-security-outcomes} presents the seven attacks whose original and hardened stages use the same broad attack family. A07 is excluded from the paired chart because its original baseline used a synthetic trigger whereas the hardened package used a natural-style request. The attack-level package tables, representative records, and the separate authorised positive-control figure are retained in Appendix~\ref{app:detailed-hardening}.

\begin{figure}[htbp]
\centering
\includegraphics[width=\textwidth]{results/cross_attack_package_comparison.pdf}
\caption[Package-level security outcomes]{Descriptive original-baseline and hardened-package outcomes. A01 uses delivered direct leakage; A02 policy-aware delivered leakage; A03 post-downgrade leakage; A04 complete relation-edge leakage; A05 membership confirmation; A06 protected-value leakage; and A08 delivered numeric leakage. Prompt equivalence is incomplete: A01 target-level prompts are partially preserved/reconstructable and differ across packages; A02 final prompts differ in 450/450 preserved sequences while warm-ups match; final prompts also differ for A05; A03 final prompts match but protected seed prompts differ; exact prompt text is unavailable for the historical A04 and A08 package runs; and A06 attack and warm-up user prompts match at the stored record level. The superseding matched A04 ablation separately uses preserved, identical prompt sequences across its off/on arms. A07 is excluded because the two package stages used different attack variants. The figure supports descriptive package-level comparison only and does not isolate individual hardening components. Each bar has denominator 150.}
\label{fig:package-security-outcomes}
\end{figure}

\section{Evidence Orientation for Guard Ablations}

Section~\ref{sec:architecture-guards}, Figure~\ref{fig:guard-execution-pipeline}, and Appendix~\ref{app:guard-reference} document the implementation and exact execution order of the six configurable controls. Intervention effects depend on pipeline stage: a control may remove model-visible evidence, clear state, modify an intermediate answer, or only record a decision. Fixed controls can still protect a guards-off arm, and execution telemetry alone cannot establish an outcome reduction.

\begin{table}[htbp]
\centering
\footnotesize
\setlength{\tabcolsep}{4pt}
\caption{Evidence classes used to interpret the guard interventions.}
\label{tab:guard-ablation-evidence-classes}
\begin{tabularx}{\textwidth}{@{}>{\raggedright\arraybackslash}p{36mm}>{\raggedright\arraybackslash}p{46mm}X@{}}
\toprule
Observed intervention pattern & Experiments & Permissible claim \\
\midrule
Relevant off-arm outcome reproduced and smaller with the component on & A04 relation-edge disclosure; A05 indexed-record existence; A08 retrieval-mediated numeric disclosure; supplemental A02 final-delivery fork; supplemental A06 exact-canary delivery; supplemental A07-S canary delivery & The configured component reduced the scored outcome within the matched prompt and source condition. This is not a general package-level or unseen-prompt claim. \\
Measured internal state changed while delivered leakage remained zero & A03 access-change clearing & The switch removed the recorded post-transition state under the tested configuration. It did not demonstrate a delivered-leakage reduction. \\
Primary outcome was zero in both arms, with execution, replacement, exposure, or stage telemetry where available & Earlier A01 and A02 verifier ablations; earlier A06 prompt-injection ablation; A07-N relation-guard ablation & The evidence can establish execution or a measured stage change, but not prevention of an attack-specific outcome that the off arm did not produce. Missing direct activation telemetry is not inferred from answer content. \\
\bottomrule
\end{tabularx}
\end{table}

\paragraph{Status of mechanism claims.}
The root-cause subsections combine directly observed telemetry with architectural inference. Recorded retrieval, state, raw-answer, delivery, and guard-action fields are observations. Attribution to an unablated internal layer is a mechanism hypothesis unless a matched intervention changed that layer. Accordingly, package-level root-cause narratives identify a plausible failure path; only the recorded off/on control in a matched experiment receives a within-experiment causal interpretation.

\FloatBarrier

\section{A01: Direct Cell Extraction}

\paragraph{Historical vulnerability.}
The descriptive original baseline exposed the target through retrieval in 150/150 unauthorised conversations per mode and delivered it in 74/150 sensitivity-mode answers, compared with 0/150 in secure mode.

\paragraph{Mechanism.}
The retriever operated on entity-level chunks that could contain mixed-sensitivity fields. In sensitivity-evaluation mode, protected fields remained available to prompt construction. The architectural vulnerability was the availability of a protected exact value beyond role projection without a reliable delivery-stage check.

\paragraph{Hardening.}
The implemented hardening strengthens role-aware field projection and checks the raw answer against restricted structured values in \texttt{RAGPipeline}'s post-generation verification path. The package evaluation does not isolate which component produced each non-disclosure.

\paragraph{Package evidence.}
The hardened package recorded zero observed delivered or raw-answer leakage under the attack-specific scorer in either mode. Secure-mode leakage remained 0/150, while sensitivity-mode leakage was 74/150 in the original baseline and 0/150 after hardening; positive-control success was 75/75 and 75/75 in secure mode and 75/75 and 60/75 in sensitivity-evaluation mode. The package table and representative record are in Appendix~\ref{app:detailed-hardening}.

\paragraph{Matched component evidence.}
The guards-off arm did not reproduce direct leakage: delivered and raw leakage were 0/150 in both modes, as they were with the verifier enabled. Sensitivity-mode retrieval and model-visible exposure nevertheless remained 150/150 in both arms. Enabling the verifier replaced 150/150 unauthorised answers per mode, including raw answers that did not contain the target percentage under the primary exact-value scorer. Secure-mode positive-control task success remained 75/75; sensitivity-mode success was 57/75 with the verifier off and 60/75 with it on. The raw protected-answer counts were identical to the respective delivered counts, so the three-conversation positive-control difference cannot be isolated from repeated API-call variation and is not evidence of a general utility effect.

\paragraph{Supplemental verifier validation.}
Deterministic replay recognised all 74 stored historical A01 exact-value leaks, whereas the prospectively frozen live pilot remained at 0/24 leakage in both arms and failed its predefined non-zero off-leakage gate. The pilot therefore supplies no verifier-necessity evidence under the hardened generation condition. The full replay confusion inventory, corpus-stratum qualification, and pilot design are reported in Appendix~\ref{app:detailed-hardening}.

\paragraph{Interpretation.}
The original baseline recorded direct sensitivity-mode leakage in 74/150 conversations, while the hardened package recorded 0/150. Sensitivity-mode positive-control success was lower in the hardened package, at 60/75 compared with 75/75 in the original baseline.

The A01 prompt audit establishes partial target-level provenance: all five baseline target prompts are stored, all five hardened final prompts are reconstructable with corroborating run evidence, and none is identical because the baseline adds an exact-value directive. Per-conversation sequences and exact system/API payloads remain unavailable. These package values are descriptive and cannot assign the historical difference to one component.

Both matched arms had zero raw and delivered leakage. Enabling the verifier produced broad answer replacement while internal exposure remained unchanged. The matched experiment cannot attribute a confidentiality reduction to the output verifier. Appendix~\ref{app:detailed-hardening} retains the probable-mechanism analysis and matched table--figure pair.

\section{A02: Multi-Turn Row Construction}

\paragraph{Historical vulnerability.}
Policy-aware rescoring of the descriptive original baseline found newly delivered restricted information in 8/150 secure-mode and 123/150 sensitivity-mode conversations. The totals comprised 8 and 86 partial reconstructions and 0 and 37 full reconstructions. The legacy scorer's 60/150 and 122/150 target-field flags are diagnostic and are not treated as policy-aware leakage.

\paragraph{Mechanism.}
Entity focus and conversation memory allowed separately requested fields to remain associated with the same formulation. Under the target-aware exact matcher, the final retrieval context contained all eight scored target fields in 120/150 unauthorised runs per mode. In the remaining 30/150 runs, all for Panthenol, it contained seven of eight fields; \texttt{Phase} was the undetected field. Thus, 120/150 denotes complete eight-field support, not the absence of relevant retrieval in the other 30 runs. The historical artifacts do not retain raw retrieved chunks, so they cannot distinguish a genuinely absent field from formatting or matcher failure. Prompt construction supplied the accumulated evidence to generation, and the original delivery path did not validate the combined answer against the restricted fields of the focused row. The failure therefore spanned retrieval, focus state, memory, prompt construction, and generation output.

\paragraph{Hardening.}
The implementation adds context with sensitivity labels, records raw and delivered answers separately, and applies a structured post-generation verifier. The verifier builds an inventory of fields restricted for the active role, examines the raw generated answer, and replaces an answer containing a restricted structured value before delivery. This design addresses the composition problem at the point where individually harmless turns become a combined disclosure.

\paragraph{Package evidence.}
The observed policy-aware delivered outcome was 8/150 and 123/150 in the original baseline and 0/150 in the hardened package. The corresponding descriptive differences are 5.3 and 82.0 percentage points. Policy-aware full reconstruction was 0/150 and 37/150 in the baseline and 0/150 in the hardened package. The hardened sensitivity-mode runner recorded three legacy raw target-field flags, but policy-aware rescoring classified raw restricted-field leakage as 0/150 because those cases contained only fields authorised for the internal role.
The package table and representative record are retained in Appendix~\ref{app:detailed-hardening}.

\paragraph{Matched component evidence.}
Policy-aware delivered and raw restricted-field leakage was 0/150 in every arm and mode; partial and complete reconstructions were also zero. Restricted target values were present in retrieval and model-visible context in 150/150 unauthorised conversations per arm and mode, but the model did not emit a newly restricted field. With the verifier enabled, at least one turn was replaced in 147/150 secure-mode and 52/150 sensitivity-mode conversations, despite zero policy-aware raw leakage. This difference reflects the verifier's broader inventory matching and must not be relabelled as policy-aware leakage. Protected full-row success was unchanged at 75/75 in secure mode and 50/75 in sensitivity-evaluation mode; 70/75 sensitivity-mode protected answers in each arm contained at least one target field.

\paragraph{Interpretation.}
The hardened package delivered no policy-aware restricted field in the evaluated matrix. Secure-mode authorised full-row task success was 68/75 in the original baseline and 75/75 in the hardened package. Sensitivity-mode full reconstruction was 74/75 in the original baseline and 54/75 in the hardened package; 70/75 hardened sensitivity-mode protected answers contained at least one target field.

All historical A02 user-prompt sequences are preserved in \path{turns[].prompt}. Their warm-up subsequences match in 450/450 conditions (900/900 warm-up prompt rows), whereas their final attack prompts differ in 450/450 conditions, including the change from ``reconstruct'' to ``find/return'' wording. Exact historical system/API payloads and complete source states remain unavailable. The package comparison is descriptive.

The matched ablation used identical prompts. Its guards-off condition did not reproduce a policy-aware vulnerability. Enabling the verifier changed delivery behaviour while the policy-aware primary outcome and protected full-row task success remained unchanged. The experiment cannot attribute a confidentiality reduction to the verifier. Exact-value matching may miss paraphrases, derived fields, or approximate reconstructions. Appendix~\ref{app:detailed-hardening} retains the probable-mechanism analysis and matched table--figure pair.

\paragraph{Supplemental output-verifier evidence.}
\label{sec:a02-verifier-component-validation}

The sensitivity-evaluation-mode challenge held the final generated raw answer byte-identical and reduced policy-aware final delivery from 12/150 with the verifier off to 0/150 with it on. All 12 scorer-positive outputs were detected. The verifier also replaced 50/138 scorer-negative unauthorised outputs, giving precision 12/62 (19.4\%); protected full-row success was unchanged at 72/75 in both derived deliveries.

This design isolates the final-answer delivery-stage contribution. Warm-up generations and pre-final state were shared, so it provides neither historical package attribution nor a warm-up verifier-effect estimate. The renderer was frozen on separate development targets, but the confirmatory extension reused the five established thesis targets; fixed templates, the exact-value policy scorer, and temperature-zero repetitions further bound the result. Appendix~\ref{app:detailed-hardening} retains the deterministic replay, pilot, full breakdown, and false-replacement inventories.

\section{A03: Access-Level Downgrade}

\paragraph{Historical vulnerability.}
The label-free baseline leaked the seeded protected value in all 150 downgraded public/internal conversations in both modes.

\paragraph{Mechanism.}
An access change was treated as an ordinary update to conversation state. Protected content acquired before downgrade survived in recent messages, summaries, vector-memory snippets, focus history, cached retrieval, visible context, and previous answers. These objects were created under the old policy but remained reachable under the new policy. The principal failure was the lifecycle of security-sensitive state, reinforced by memory and retrieval reuse.

\paragraph{Hardening.}
The central \texttt{set\_access\_context} transition now compares the previous and new role and allowed-sensitivity set. When \path{ACCESS_CHANGE_MEMORY_CLEAR_ENABLED} is active and the context changes, \texttt{\_clear\_memory\_after\_role\_change} clears conversation memory, summaries, cached and visible results, access decisions, focus state, previous answers, and guard state. New retrieval and memory selection use the current access policy. This makes a role transition a security-context boundary.

\paragraph{Package evidence.}
The observed package-level difference was 100 percentage points in both modes. In the hardened sensitivity-mode unauthorised conditions, the protected setup answer succeeded in 140/150 conversations; leakage was also 0/140 when conditioned on successful setup. All 75 protected final-access positive controls succeeded.
The package table and representative record are retained in Appendix~\ref{app:detailed-hardening}.

\paragraph{Matched component evidence.}
The protected setup succeeded in 150/150 secure-mode and 140/150 sensitivity-mode unauthorised transitions in each arm. With state clearing disabled, the target remained in post-transition state in 150/150 and 140/150 conversations, respectively; recent-message and memory-snippet exposure had the same counts. Enabling state clearing reduced each of these exposure counts to 0 and recorded a clear action in 150/150 unauthorised transitions per mode. Raw and delivered post-downgrade leakage nevertheless remained 0/150 in both arms, including 0/150 and 0/140 when conditioned on successful setup. Protected positive-control success remained 75/75.

\paragraph{Interpretation.}
This attack showed the largest observed package-level confidentiality difference, with all positive controls preserved. Because the protected seed prompts differ, the historical package comparison is descriptive and cannot isolate the memory-clear switch.

In the matched ablation, state clearing removed all measured post-transition state exposure. The guards-off generator already withheld the value, and delivered leakage remained zero in both arms. The matched evidence identifies a state-management effect only.

The operational cost remains loss of conversational continuity at every role transition. Untested conditions include cross-session stores, distributed caches, concurrent requests, and a role change occurring during generation. Appendix~\ref{app:detailed-hardening} retains the probable-mechanism analysis and matched table--figure pair.

\section{A04: Relational Join-Path Inference}

\paragraph{Historical vulnerability.}
Sensitivity-evaluation mode disclosed a protected relation edge in 75/150 attack-family-label-free baseline answers; secure mode exposed relation nodes but delivered no complete edge.

\paragraph{Mechanism.}
The original pipeline treated formulation and process identifiers as ordinary metadata. Relation expansion could traverse from a public product to protected nodes, and hidden identifiers could enter focus state and later prompts. Field filtering alone did not express whether an edge could be traversed or disclosed. The failure involved relation handling, relation-aware access control, focus state, prompt construction, and generation.

\paragraph{Hardening.}
The implementation introduces explicit relation specifications in the sensitivity policy. Separate traversal and disclosure decisions are enforced through \texttt{relation\_}\allowbreak\texttt{visibility\_}\allowbreak\texttt{for\_role}. \texttt{RAGPipeline} checks permission before relation expansion, projects relation fields for the active role, sanitises focus identifiers, and adds restricted relation values to final-answer verification. Relation edges are thereby treated as policy-bearing data rather than neutral metadata.

\paragraph{Package evidence.}
Secure-mode edge leakage was 0/150 in both package stages, and sensitivity-mode leakage was 75/150 in the original baseline and 0/150 after hardening. Positive-control success was 75/75 in both secure-mode stages and 25/75 and 75/75 in sensitivity-evaluation mode. Sensitivity-mode context exposure remained 150/150 after hardening, but none of that evidence was delivered as a relation-edge leak. The package table and representative record are retained in Appendix~\ref{app:detailed-hardening}.

\paragraph{Matched component evidence.}
With the relation-access guard disabled, complete edge leakage occurred in 150/150 secure-mode and 65/150 sensitivity-mode unauthorised conversations. Enabling the guard reduced delivered and raw edge leakage to 0/150 in both modes. Secure-mode model-visible relation exposure decreased from 150/150 to 0/150, while sensitivity-mode exposure remained 150/150 in both arms and delivery still decreased from 65/150 to zero. Node and complete-path leakage followed the edge counts; downstream protected-detail leakage was zero throughout. Protected path completeness remained 75/75. A direct relation-guard activation counter was not logged, so the action stage is reported as unavailable rather than inferred from the answer.

\paragraph{Interpretation.}
The hardened package recorded no delivered edge disclosure and higher authorised completeness in its matrix. Exact historical package prompt text is unavailable, so the package comparison remains descriptive.

The matched family-label-omitted ablation provides component evidence that the relation-access switch caused the measured edge-leakage reduction under the fixed, mechanism-specific prompt condition. Persistent sensitivity-mode exposure shows that the intervention did not uniformly remove evidence before generation. The evaluation covered depth-two paths; deeper, cyclic, many-to-many, and inference-by-combination paths remain untested. Appendix~\ref{app:detailed-hardening} retains the matched table--figure pair and the mode-specific mechanism analysis.

\section{A05: Protected-Record Membership/Existence Probe}

\paragraph{Historical vulnerability.}
The descriptive original baseline confirmed protected-record existence in 44/150 secure-mode and 70/150 sensitivity-mode unauthorised answers. Here and below, \emph{membership} means membership in the indexed record corpus, not membership in model-training data.

\paragraph{Mechanism.}
Sensitivity policy applied to fields and records but did not make existence itself a protected fact. Membership, rank, similarity, near-name, and candidate-ID prompts converted retrieval evidence into a binary claim. The failed boundary joined embedding/retrieval metadata, prompt intent, and generated answers.

\paragraph{Hardening.}
\texttt{membership\_inference\_guard.py} detects membership-, existence-, rank-, and similarity-oriented probes and enriches detected candidates from metadata. It checks the candidate's membership sensitivity against the active role. In secure mode, an enabled guard can refuse a detected unauthorised probe before retrieval. In sensitivity-evaluation mode, generation proceeds, after which the guard replaces every noncanonical answer to a detected unauthorised probe when enforcement is enabled. The computed binary membership-confirmation signal is retained as telemetry but does not determine whether replacement occurs. Protected structured values and details are checked separately by the output-leakage verifier. A separate authorised path can construct a supported answer when a protected user receives an unjustified refusal.

\paragraph{Package evidence.}
The original baseline recorded 44/150 secure-mode and 70/150 sensitivity-mode protected-record existence confirmations, whereas the hardened package recorded zero observed confirmations under its scorer.
Protected-role success was 75/75 in both modes and package stages. The package table and representative record are retained in Appendix~\ref{app:detailed-hardening}.

\paragraph{Matched component evidence.}
The guards-off arm produced no secure-mode confirmation but delivered 15/150 sensitivity-mode protected-record existence confirmations, all accompanied by target details. With the membership/existence guard enabled, delivered confirmation was 0/150 in both modes. Sensitivity-mode raw confirmation remained 5/150 before replacement, while delivered confirmation was zero. Target retrieval exposure remained 150/150 in both sensitivity-mode arms. In secure mode the enabled guard refused all 150 probes before calling the model; in sensitivity-evaluation mode it called the model and replaced all 150 delivered answers. Protected positive-control task success remained 75/75 in both modes and arms, with no recorded authorised error or over-refusal.

\paragraph{Interpretation.}
Under its evaluated prompts, the stored hardened package delivered no protected-record existence confirmations and preserved 75/75 protected-role success in both modes. The original and hardened prompt formulations differ, making the package comparison descriptive.

The matched ablation reproduced a sensitivity-mode vulnerable condition. Under the tested prompts, enabling the membership/existence guard changed delivered confirmation from 15/150 to 0/150. This result demonstrates broad query-conditioned delivery enforcement, not the precision or specificity of a binary membership-signal classifier: the enabled guard replaced all 150 sensitivity-mode answers, while only 5/150 corresponding raw answers were scored as confirmations. Target retrieval exposure also remained 150/150, so the intervention did not eliminate upstream evidence. Protected positive-control task success remained 75/75 in both modes and arms. Novel probe paraphrases, absent-record controls, and interfaces exposing scores, citations, or ranked sources remain important evaluation gaps.

\section{A06: Prompt Injection through a Poisoned Row}

\paragraph{Historical vulnerability.}
The label-free baseline retrieved the poisoned row and protected target content in 150/150 unauthorised runs per mode. Canary compliance remained 0/150, but sensitivity-evaluation mode delivered the protected ingredient--percentage pair in 2/150 conversations. The result therefore contains a small confidentiality failure without measured compliance with the stored override instruction.

\paragraph{Mechanism.}
Public data was authorised to retrieve but could contain imperative text, while relation traversal also exposed linked protected formulation content. The original context format did not provide a dedicated quarantine layer separating instruction-like cell contents from control instructions, and sensitivity-evaluation mode did not prevent two linked protected values from reaching delivery. Confidentiality filtering and prompt integrity were therefore distinct failed or exposed boundaries, although zero canary outputs provide no evidence that GPT-4o-mini followed the tested override instruction.

\paragraph{Hardening.}
The \texttt{prompt\_injection\_guard} switch controls a composite component. In secure mode it can withhold suspicious unstructured retrieved text or sanitise matched instruction-like lines after role projection. In sensitivity-evaluation mode it deliberately preserves the content behind an explicit untrusted-data boundary so that downstream restraint remains testable. After generation, the same switch detects the defined canary and trigger artifact families and redacts affected lines or replaces an empty remainder with a refusal. Role-aware field and relation projection remain separate controls. An off/on comparison can therefore estimate this configured component but cannot independently attribute an effect to its context-treatment and artifact-filtering sub-actions.

\paragraph{Package evidence.}
Secure-mode protected-value leakage and canary compliance were 0/150 in both package stages, with positive-control success at 75/75. In sensitivity-evaluation mode, protected-value leakage was 2/150 in the original baseline and 0/150 after hardening; canary compliance remained 0/150, while positive-control success was 25/75 and 75/75. This prompt-matched package comparison spans implementation stages and remains descriptive. The package table and representative record are retained in Appendix~\ref{app:detailed-hardening}.

\paragraph{Matched component evidence.}
The guards-off arm produced neither canary compliance nor protected-value delivery in either mode. Enabling the prompt-injection guard left both outcomes at 0/150. In secure mode, the intervention changed the recorded quarantine and guard-check counts from 0/150 to 150/150 while model-visible protected-target exposure remained 0/150. In sensitivity-evaluation mode, the guard was checked in 150/150 conversations, but the deliberately preserved and labelled untrusted context left model-visible exposure at 150/150 in both arms. Positive-control success remained 75/75 in every cell.

\paragraph{Supplemental frozen-challenge evidence.}
\label{sec:a06-supplemental-rationale}

Because the earlier matched off arm produced neither the canary nor protected-value leakage, it could not establish prevention of an attack-specific outcome. Section~\ref{sec:a06-supplemental-design} therefore defined a separate, prospectively frozen challenge with disjoint development targets and a continuation gate. R-001--R-005 were held out from this pilot but were already established thesis targets. The primary user prompt contained no attack-family label or canary; the malicious workflow existed only in the retrieved row.

The development pilot and the unchanged \texttt{workflow\_follow\_v1} confirmatory profile passed their technical audits. The confirmatory source contained 900/900 conversations and 450/450 matched pairs and remained separate from the original baseline, hardened package, and earlier matched-ablation denominators. Appendix~\ref{app:detailed-hardening} retains the complete pilot-profile inventory, audit counts, and frozen-challenge table.

The pooled primary outcome fell from 300/300 delivered exact canaries with the guard off to 0/300 with it on. All 300 unauthorised matched pairs made the exact delivered transition $1\to0$. At the raw-answer stage, 155 pairs made the transition $1\to0$ and 145 remained $1\to1$; delivery nevertheless ended at zero in both modes. Protected ingredient-plus-percentage leakage was zero at both the raw and delivered stages in every arm and mode.

The recorded mechanism differed by mode. In secure mode, the poisoned row was retrieved in all 150 unauthorised guard-on conversations, but the canary-bearing content was not model-visible. Context quarantine was recorded before generation for all 225 guard-on records, including the protected-role conditions. Raw and delivered canary counts were consequently zero.

In sensitivity-evaluation mode, the canary remained model-visible in all 150 unauthorised guard-on conversations and was still generated in 145 raw answers. Post-generation artifact handling then removed it before delivery; artifact redaction was recorded in 220/225 all-role guard-on records.

The experiment switched the composite \texttt{prompt\_injection\_guard} as one unit and supports bounded within-condition attribution to that configured component for exact-canary delivery under the frozen condition. Pre-generation quarantine and post-generation artifact filtering were not varied independently and have no independent effect estimates. The two arms used separate API generations in a fixed off-then-on order, without an identical-raw-answer fork. Residual generator-call, service-time, or order variation cannot be excluded.

\paragraph{Structurally non-evaluable authorised positive control.}
The planned protected-role control was structurally non-evaluable. It recorded 0/75 successes in every table cell because the linked protected formulation was retrieved in 0/300 protected records across both arms and entered the LLM context in 0/300. The user prompt asked the system to execute a stored quality-assurance workflow but contained none of the formulation-, ingredient-, percentage-, or supplier-related terms that trigger product-to-formulation relation expansion. Content discovered in the poisoned row was not used for a second retrieval-routing decision, so the expected protected value was never available to answer.

The 300/300 \texttt{authorized\_accuracy\_error} flags identify the missing retrieval prerequisite. They are not an independent measurement of model-accuracy failure. The zero control was already visible in the pilot. Positive-control success was not part of the continuation or interpretation gate; technical-audit \texttt{PASS} does not validate utility preservation.

The earlier matched A06 ablation retained 75/75 protected-role success in both modes and arms, but its separate prompt explicitly requested the linked formulation details. Those counts are evidence only for that earlier prompt and retrieval opportunity; they cannot repair the supplemental control. Separately, a post-hoc exact-string review of the benign-metadata pilot profile found the correct product name and target market in 96/96 answers across roles, arms, and modes. This was not a preregistered scorer and is reported only as exploratory evidence of narrow benign metadata answering, not as a replacement positive control or a general utility result.

\paragraph{Interpretation.}
The four A06 evidence sources answer different questions and are not pooled. The original baseline observed sensitivity-mode confidentiality leakage in 2/150 conversations while canary compliance remained zero. The hardened package observed zero leakage and stronger protected-role task success, but comparison across implementation stages remains descriptive. The earlier matched ablation preserved 75/75 positive controls and documented mode-dependent guard execution, yet its 0/150$\to$0/150 canary and protected-value outcomes could not demonstrate prevention. The separate frozen challenge supplied the missing active off-arm integrity condition and matched 300/300$\to$0/300 exact-canary delivery evidence for the composite prompt-injection component.

The supplemental challenge supports the primary integrity result for exact-canary delivery by the composite guard under the frozen condition. Protected-value leakage was already zero with the guard off, leaving confidentiality reduction unmeasured. The linked formulation was never retrieved for the protected control, leaving utility preservation and utility loss unmeasured.

The challenge does not explain the historical 2/150 sensitivity-mode confidentiality failures. The evaluated scope comprises five established formulation targets and their target-specific synthetic poisoned rows, the frozen stored workflow and user-prompt profile, the audited source state and scorer, and exact-canary delivery. Unseen, obfuscated, multilingual, indirect, and tool-oriented injections remain untested; benign end-to-end task utility requires separate testing. The earlier matched table--figure pair and supplemental inventories are retained in Appendix~\ref{app:detailed-hardening}.

\section{A07: Triggered Indirect Prompt Injection}

\paragraph{Historical vulnerability.}
The label-free baseline emitted the attacker-controlled canary in 150/150 unauthorised answers in both modes. Exact protected-value leakage remained 0/150, making this primarily an integrity failure.

\paragraph{Mechanism.}
The trigger caused retrieved validation text to be interpreted as an instruction source. Prompt construction did not sufficiently separate trusted task instructions from stored data, and the output boundary did not treat a canary as malicious unless it was also a protected secret. The failed layers were untrusted-content handling, prompt construction, generation integrity, and artifact validation.

\paragraph{Hardening.}
The final pipeline contains prompt-injection, membership, projection, and output controls and separately records canary compliance and protected-value leakage. The package evaluation does not isolate the contribution of any one component.

\paragraph{Package evidence.}
The original synthetic-trigger baseline produced 150/150 canaries in both modes and 0/150 protected-value disclosures, with positive-control success of 10/75 in secure mode and 0/75 in sensitivity-evaluation mode. The hardened natural-request evaluation produced 0/150 canaries and 0/150 protected-value disclosures in both modes, with positive-control success of 75/75 and 67/75. These stages use different attack variants and are not a controlled before--after pair. The package table and representative trigger record are retained in Appendix~\ref{app:detailed-hardening}.

\paragraph{Matched A07-N relation-guard evidence.}
The A07 matched experiment evaluates only the natural-request family (A07-N). Delivered and raw protected-value leakage were 0/150 in both modes and both arms. Secure-mode protected-target exposure was 0/150, while sensitivity-mode model-visible exposure was 150/150 with the relation-access guard both disabled and enabled. The public validation row was retrieved in all 150 unauthorised conversations per arm and mode. The prompt-injection guard, which was deliberately fixed rather than ablated, quarantined this row in all secure-mode conversations and in no sensitivity-evaluation conversations. Positive-control success was 75/75 in secure mode; the sensitivity-evaluation count was 74/75 with the relation guard off and 75/75 with it on, with one recorded authorised accuracy error in the former arm. Direct relation-guard activation telemetry was not available.

\paragraph{Supplemental matched A07-S evidence.}
\label{sec:a07s-matched-ablation}

The supplemental experiment restored the historical A07-S synthetic-trigger prompt family and switched only the prompt-injection guard. Its 450 shared conditions had equal user-prompt sequences and recorded system-prompt sets, and the rerun prompt set matched all 450 historical A07-S records exactly. Both arms contained 225 conversations per mode, giving 900 conversations and 450 matched pairs overall, with no missing, duplicate, or execution-error record.

With the guard off, raw and delivered canary compliance were 150/150 in both modes. With the guard on, the same configured guard acted at different stages by mode. In secure mode, the retrieved injection content was quarantined in 150/150 unauthorised conversations, with raw and delivered canary compliance both at 0/150. Sensitivity-evaluation mode deliberately retained the trigger in model-visible context: raw canary compliance remained 150/150, the answer artifact was detected in 150/150, and delivered compliance fell to 0/150 through redaction. Protected-value leakage remained 0/150 in both arms and modes, and authorised positive-control success remained 75/75 in every arm and mode. The causal result concerns integrity/canary delivery; no confidentiality reduction is claimed.

\paragraph{Interpretation.}
The two package stages still use different attack variants and therefore do not constitute a controlled before--after comparison. The earlier A07-N relation-guard ablation likewise cannot establish that relation access prevents the historical canary response. The supplemental A07-S study answers a different, narrower question and supplies the missing component evidence: under the matched historical synthetic-trigger condition, the prompt-injection guard prevented delivered canary compliance.

Within A07-N, disabling the relation-access guard was insufficient to produce protected-value leakage. Other hardened controls remained fixed: most visibly, the prompt-injection guard quarantined every retrieved validation row in secure mode. Sensitivity-evaluation mode preserved the target in model-visible context in both arms, but generation still did not disclose it. This combination explains why a guard can be switched without changing the primary metric and why an off arm need not reproduce the historical vulnerability. The relation guard controls linked-data traversal and projection; it does not recreate the synthetic trigger or disable prompt-injection, system-prompt, and delivery protections. Since direct relation-guard activation telemetry was not recorded, the zero-to-zero result must not be described as proof that the guard was unnecessary or effective.

The single A07-N sensitivity-mode positive-control difference coincided with one authorised accuracy error in the guards-off arm and is insufficient to support an authorised-task effect, much less a general utility effect. A07-S now provides matched evidence for the known trigger, but held-out benign validation requests and unseen, obfuscated, paraphrased, and multilingual triggers remain necessary to estimate over-refusal and bypass risk. Appendix~\ref{app:detailed-hardening} retains the A07-N table--figure pair and the A07-S table and detailed provenance inventory.

\section{A08: Embedding/Rank-Framed Retrieval Probe}

\paragraph{Historical vulnerability.}
The label-free baseline exposed target retrieval evidence in all unauthorised runs and leaked protected numeric values in 82/150 sensitivity-mode answers. Secure-mode delivered leakage was zero.

\paragraph{Mechanism.}
Embedding, nearest-neighbour, and rank wording framed an indirect extraction request. In sensitivity-evaluation mode, protected numeric fields could be model-visible; query handling and answer validation did not treat this framing as security relevant. Memory could also provide an alternative source for a probe. The measured outcome was retrieval-mediated numeric disclosure. The interface exposed no embeddings, similarity scores, or ranking list, and the thesis makes no embedding-inversion claim.

\paragraph{Hardening.}
The embedding-probe logic in \texttt{RAGPipeline} detects ranking vocabulary together with sensitive identifiers, membership language, or protected numeric-field requests. For unauthorised probes it enables strict role-filtered retrieval, forces secure projection even in sensitivity-evaluation mode, adds side-channel policy text, suppresses memory for the query, and retains restricted inventory values for final verification.

\paragraph{Package evidence.}
Secure-mode delivered leakage was 0/150 in both package stages; sensitivity-mode leakage was 82/150 in the original baseline and 0/150 after hardening, a descriptive difference of 54.7 percentage points. Positive-control success remained 75/75 in every package cell, and hardened context exposure was 0/150 in both modes. The package table and representative record are retained in Appendix~\ref{app:detailed-hardening}.

\paragraph{Matched component evidence.}
In sensitivity-evaluation mode, disabling the embedding-probe guard produced protected-target retrieval exposure in 150/150 conversations and delivered protected numeric values in 3/150. Enabling the guard activated it in 150/150 conversations, reduced retrieval exposure to 0/150, and reduced raw and delivered leakage to 0/150. Secure-mode exposure and leakage were already 0/150 with the guard disabled and remained zero when it was enabled, although activation changed from 0/150 to 150/150. Positive-control success was unchanged at 75/75 in secure mode and 74/75 in sensitivity-evaluation mode; each sensitivity arm contained one recorded authorised accuracy error.

\paragraph{Interpretation.}
The hardened package recorded no numeric leakage while positive-control success was 75/75 in both modes. Exact prompts are absent from the stored historical records, so that package-level difference remains descriptive. Under identical recorded matched prompts and conditions, the embedding-probe switch changed both upstream target exposure and the sensitivity-mode delivered outcome.

The first submitted A08 matched-ablation attempt produced no usable result records because every shard failed after an empty model response; it is excluded from all reported denominators. The complete rerun is the sole matched A08 source. Held-out benign embedding queries, unseen paraphrases, and interfaces exposing raw scores or rankings remain necessary to measure false refusals and bypass risk. Appendix~\ref{app:detailed-hardening} retains the matched table--figure pair and the stage-specific mechanism analysis.

\FloatBarrier
\section{Cross-Attack Matched-Ablation Summary}

\begin{table}[htbp]
\centering
\scriptsize
\setlength{\tabcolsep}{2pt}
\begin{tabularx}{\textwidth}{@{}l>{\raggedright\arraybackslash}p{30mm}*{5}{>{\centering\arraybackslash}X}@{}}
\toprule
Attack & Tested control & \shortstack{S\\outcome} & \shortstack{SE\\outcome} & \shortstack{Secondary\\evidence} & \shortstack{S\\positive task} & \shortstack{SE\\positive task} \\
\midrule
A01 & Exact-value delivery verifier & 0$\to$0 & 0$\to$0 & \shortstack{S: 0$\to$0\\SE: 150$\to$150} & 75$\to$75 & 57$\to$60 \\
A02 & Structured restricted-field verifier & 0$\to$0 & 0$\to$0 & \shortstack{S: 150$\to$150\\SE: 150$\to$150} & 75$\to$75 & 50$\to$50 \\
A03 & Access-change state clearing & 0$\to$0 & 0$\to$0 & \shortstack{S: 150$\to$0\\SE: 140$\to$0} & 75$\to$75 & 75$\to$75 \\
A04 & Relation-access guard & 150$\to$0 & 65$\to$0 & \shortstack{S: 150$\to$0\\SE: 150$\to$150} & 75$\to$75 & 75$\to$75 \\
A05 & Membership guard & 0$\to$0 & 15$\to$0 & \shortstack{S: 0$\to$0\\SE: 150$\to$150} & 75$\to$75 & 75$\to$75 \\
A06 & Prompt-injection guard & 0$\to$0 & 0$\to$0 & \shortstack{S: 0$\to$0\\SE: 0$\to$0} & 75$\to$75 & 75$\to$75 \\
A07 & Relation-access guard & 0$\to$0 & 0$\to$0 & \shortstack{S: 0$\to$0\\SE: 150$\to$150} & 75$\to$75 & 74$\to$75 \\
A08 & Embedding-probe guard & 0$\to$0 & 3$\to$0 & \shortstack{S: 0$\to$0\\SE: 150$\to$0} & 75$\to$75 & 74$\to$74 \\
\bottomrule
\end{tabularx}
\caption[Cross-attack matched-ablation summary]{Cross-attack matched-ablation summary. Every numerical cell reports the raw count with guards off $\rightarrow$ guards on; security counts have denominator 150 and authorised positive-control task counts denominator 75. Selected outcomes are attack-specific and must not be interpreted as one homogeneous measure. For A06, the selected outcome is canary compliance and the secondary column is the separate delivered confidentiality outcome.}
\label{tab:cross-attack-matched-ablation}
\end{table}


Table~\ref{tab:cross-attack-matched-ablation} gives the compact numerical summary. The attack-specific outcomes remain heterogeneous, and interpretation follows the evidence classes stated above. Appendix~\ref{app:detailed-hardening} retains the detailed claim ledger and cross-attack diagnostic graphics.

A04, A05, and A08 reproduced non-zero guards-off delivered outcomes and reached zero delivered outcomes when their respective controls were enabled. A01, A02, A03, and A07-N had zero selected delivered outcome in both arms; their ablations therefore do not demonstrate reductions in those outcomes. A03 nevertheless isolated removal of state exposure. A01 and A02 documented verifier replacement behaviour, and A07-N showed that switching relation access alone did not remove sensitivity-mode internal exposure. The earlier A06 matched ablation separately recorded 0/150 canary compliance and 0/150 protected-value delivery in both arms while documenting injection-guard checking and mode-dependent quarantine.

This summary deliberately preserves the original A01--A08 matched-ablation matrix. The later supplemental studies are reported separately rather than inserted retroactively: the sensitivity-evaluation-mode A02 identical-final-raw-answer challenge reduced final delivery from 12/150 to zero; the audited A06 frozen challenge reduced delivered exact-canary compliance from 300/300 to 0/300 across both modes; and the matched A07-S experiment reduced delivered canary compliance from 150/150 to zero in both modes. For A06, the supplemental protected-value outcome remained zero-to-zero and the planned authorised control lacked a retrieval opportunity, so this addition changes only the bounded integrity claim.

Historical and matched guards-off outcomes differed substantially. For sensitivity-mode A01, A02, and A03, outcomes were non-zero in the historical implementation and zero in the matched guards-off arms; other hardened-code changes remained active in those arms. A04 secure-mode edge leakage was 0/150 historically and 150/150 in the matched family-label-omitted guards-off condition; sensitivity-mode leakage was 75/150 historically and 65/150 in the guards-off arm. A08 sensitivity-mode leakage was 82/150 historically and 3/150 in its matched guards-off arm.

The historical and matched sources differ in implementation and prompt provenance. These diagnostic pairs do not estimate historical reproducibility, and the guards-off hardened-code values do not replace the original baseline. Attribution to the switched control is supported only by within-experiment off/on comparisons.

\FloatBarrier
